\section{How the claim is to be measured}
\label{sec:method}

The protocol below is what the lane that fills Section~\ref{sec:results} is to
implement, in \texttt{hanzoai/benchmarks}, where it has not yet been published.
It is written here before any run exists so that the analysis cannot be chosen
after seeing the results.

\subsection{Corpora}

Corpus sizes are $N \in \{10^5, 10^6, 10^7, 10^8, 10^9\}$, drawn as nested
samples of one source collection, so every row present at one size is present
at every larger one and the growth of a cost counter is not confounded by a
change of content. The first pass is configured on English Wikipedia passages
(the 2022-12 dump, 35{,}167{,}920 passages) with published all-MiniLM-L6-v2
vectors~\cite{reimers2019sbert}, sampled by a hash threshold on the article
identifier, so that each size is a set of whole articles and a subset of the
next. That collection reaches $3.5 \times 10^7$ rows; the two larger sizes need
another source, which is not chosen. The first pass covers $10^5$ to $10^7$ on
upstream ClickHouse~26.4.2.10, the release \texttt{hanzoai/datastore} is forked
from~\cite{hanzodatastore}; the routed store follows. The collection, its licence
and its digest are recorded by the harness and printed in the caption of
Table~\ref{tab:cost}.

Part geometry is controlled rather than observed. Merge settings are fixed
across sizes, no run collapses the table to a single part, and the part count
and size distribution at each $N$ are reported with the results: routing is a
structure over parts, and a table with one part has nothing to route.

\subsection{Queries, ground truth, recall}

A query set is sampled once from the intended distribution and held fixed
across sizes. The first pass samples Natural Questions
(open)~\cite{kwiatkowski2019nq,lee2019orqa}, questions posed to a search engine
that Wikipedia answers: 200 from the training split to set the knobs and 1{,}000
from the validation split to report. A question is embedded with the corpus
model for the dense leg, and its content tokens, tokenised as the store's text
index tokenises the column, form the lexical leg. Exact answers are computed per size by exhaustive scoring, which
is affordable offline and is what makes recall meaningful; the ground-truth
procedure and its cost are recorded, as billion-scale benchmarks
do~\cite{simhadri2022bigann}.

Cost is read at a held recall target. For each system and each $N$, the knobs
--- graph search effort, part budget, the routing slack $\varepsilon$ --- are
set to the cheapest configuration whose mean recall@$k$ reaches $\rho$, and the
chosen values are published with the row. Reporting cost at a fixed knob
instead would let one system's cost fall because its answers did, which is the
failure mode that recall-against-throughput curves exist to
prevent~\cite{aumuller2020annb}.

\subsection{Counters}

Each query records, from the store's own profile events and from
instrumentation:

\begin{itemize}[leftmargin=*,itemsep=1pt]
\item \texttt{SelectedParts}, \texttt{SelectedPartsTotal},
  \texttt{SelectedMarks} and \texttt{SelectedMarksTotal} --- what index
  analysis kept and what it started from;
\item \texttt{ReadCompressedBytes} and \texttt{OSReadBytes} --- bytes from the
  store and bytes from the device;
\item routing summaries examined, posting blocks decoded, posting blocks
  skipped by a bound;
\item recall@$k$ against the exact answer, and latency.
\end{itemize} Cold and warm runs are reported separately, cold meaning the page
cache was dropped before the query. Latency is reported at the median and the
99th percentile: a design justified by an expectation has to show its tail, and
the reasons are the familiar ones~\cite{dean2013tail}.

Two derived quantities carry the argument. \emph{Cost per correct answer}
divides each counter by the number of returned rows that are in the exact
top-$k$: parts, granules and bytes per correct answer. It is the quantity to
compare across systems, because it prices recall and work in one number, and it
is the quantity the exponent is fitted to. \emph{Oracle ratio} is
$C_{\mathrm{gran}}/|G^{*}(Q)|$, which says how much more than the answer a query
had to read.

\subsection{Fitting an exponent}

For each system and for parts, granules and bytes touched per correct answer,
$C(N) = aN^{\alpha}$ is fitted by regressing $\log \mathbb{E}[C]$ on $\log N$ by
least squares over the size grid, at the held recall target. The interval on
$\alpha$ comes from a bootstrap over queries: resample the query set with
replacement, recompute the per-size means, refit~\cite{efron1993bootstrap}. The
same is done for the raw counters, and for the 99th percentile of each, since a
mean exponent below one is consistent with a tail exponent at one.

Local slopes between consecutive sizes are reported beside the global fit,
because curvature is what distinguishes a genuine exponent from an average of
two regimes. A straight line through four or five points spanning three decades
is weak evidence for a power law, and we do not claim the functional form: the
claim is about the interval on $\alpha$ under this protocol, which is a weaker
and checkable statement~\cite{clauset2009powerlaw}.

\subsection{Systems compared}

\begin{itemize}[leftmargin=*,itemsep=1pt]
\item \textbf{scan} --- exhaustive scoring with no index, the upper bound on
  cost and the reference for correctness.
\item \textbf{base} --- upstream ClickHouse 26.4: primary index, skipping
  indexes, text index for membership, one vector similarity index per part.
\item \textbf{route} --- the design of this paper.
\item \textbf{oracle} --- the granules and parts holding the exact top-$k$,
  which no system can beat and every system is measured against.
\end{itemize}

\subsection{Strata, including the ones expected to fail}

Results are reported over the whole query set and over strata: by the size of
the candidate set $\sum_{i \in A(Q)}|L_i|$, by the margin
$\theta_k - \theta_{k+1}$, by the number of basis types the query activates,
and by whether the query carries a selective predicate. Two adversarial strata
are run deliberately --- a query over an element every row carries, and a query
whose score distribution is flat enough that no region bound separates --- and
their rows are expected to show cost linear in $N$. A protocol that reported
only the mean over a benign distribution could not tell a design that routes
from one that got a convenient query set.
